\chapter{读这本书前先去掉神秘感}

\section{本章要解决的困惑}

当前 CCUSDT 策略不是一个复杂模型。它没有深度学习，没有动态 L2 路径预测，也没有复杂的最优控制器。它的 runtime 主体是一个手写状态机：

\begin{lstlisting}
L2/trade market facts
-> decision_frame row
-> TFI trigger
-> past_event / frames / spread q70 filters
-> R5 and frames cell
-> gamma sizing
-> 3x capacity ledger
-> fixed60 exit
-> entries / exits / daily / summary
\end{lstlisting}

如果只看这个链条，代码不应该难。真正让人难受的原因，是历史文档把研究假设、工程实现、运行时边界和输出审计混在一起讲。本书的目标就是把它拆回一个新人能手算的系统。

\section{全书阅读顺序}

本书不按源码目录讲，而按“一个市场事实怎样变成一笔回测交易”讲：

\begin{enumerate}[leftmargin=2em]
  \item 第一章只讲策略定义，用 L2 订单簿和成交手算，不先看工程。
  \item 第二章讲我们拿到的数据长什么样。
  \item 第三章讲数据怎样变成 \code{decision_frame}。
  \item 第四章讲 Python fast backtest 怎样逐行执行状态机。
  \item 第五章讲输出 artifacts 怎么读。
  \item 第六章讲 Python Bot 怎样在线维护同一套状态。
  \item 第七、八章讲 Rust Runner 怎样把意图变成交和日志。
  \item 后面讲 fast 和 Runner 的边界，以及新人怎样安全自己回测。
\end{enumerate}

\section{当前策略的真实复杂度}

先把核心逻辑写成伪代码：

\begin{lstlisting}
if abs(TFI) >= 1:
    direction = sign(TFI)
else:
    no entry

flat = abs(past_event_25_bps) <= 0.5
stale = frames_since_mid_change >= threshold
spread_ok = entry_cross_bps <= prior_day_q70

membership_set = trigger classes from TFI/flat/stale
base_weight = BASE_WEIGHTS[membership_set]

R5 = positive closed pnl bps / abs(negative closed pnl bps)
cell = _cell(R5 >= 1, frames_since_mid_change >= 31)
target_exposure = (base_weight + overlay) * gamma[cell]
actual_exposure = min(target_exposure, 3 - open_exposure)
\end{lstlisting}

这段没有任何神秘技巧。你真正需要理解的是每个字段从哪里来、什么时候能被策略看到、它在 L2 盘口里是什么意思。

\section{本书的硬边界}

为了不再把研究报告误读成 runtime 事实，全书反复使用下面几条边界：

\begin{lstlisting}
Runner 不读 label。
Bot 不读 date/。
Monitor 只展示，不创造事实。
Delta_bps / Energy_bps 没进入当前 CC runtime 控制。
OFI / MLOFI 不是当前 entry trigger。
\end{lstlisting}

这些话不是口号。它们决定一条数据能不能进入策略。如果一个字段必须等 entry 之后才知道，那么它只能是 label 或研究输出，不能是 runtime input。

\section{读完本书应该会什么}

读完后，你应该能做到：

\begin{itemize}[leftmargin=2em]
  \item 给一张 L2 订单簿，算 bid、ask、mid、spread、taker entry/exit 成本。
  \item 给几笔 trade，算 TFI 和方向。
  \item 给一行 \code{decision_frame}，判断是否触发 entry。
  \item 给最近 5 笔 closed outcomes，算 R5 和 cell。
  \item 给 open exposure，算 target 和 actual exposure。
  \item 给 entry/exit quote，算 fixed60 mid 和 executable PnL。
  \item 打开 Python 和 Rust 文件，知道先看哪段，不被无关工程细节淹没。
  \item 用已有数据自己跑一次 fast backtest，并知道输出结果怎么检查。
\end{itemize}

\checkpoint{本章只要记住一句话：这套系统难读，不是因为策略高级，而是因为同一个简单状态机有数据、fast、Bot、Runner、日志五种形态。}

